\documentclass[12pt,a4paper]{article}

% ==================================================
% 1. CONFIGURATION
% ==================================================

% Enable / disable title page
\newif\ifshowtitlepage
\showtitlepagetrue
% \showtitlepagefalse

% Document information
\newcommand{\worktitle}{Sprint 0 Progress Report}
\newcommand{\worksubtitle}{NAO Robot Exercise System}
\newcommand{\teamname}{NAOMI}
\newcommand{\subject}{NAO Robot Exercise System}

\newcommand{\university}{Slovak University of Technology in Bratislava}
\newcommand{\faculty}{Faculty of Informatics and Information Technologies}

\newcommand{\academicyear}{2026/2027}
\newcommand{\location}{Bratislava}
\newcommand{\documentyear}{2026/2027}

% Team leaders
\newcommand{\leaderone}{Ing. Jakub Perdek, PhD.}
\newcommand{\leadertwo}{doc. Ing. Fedor Lehocki, PhD.}

% Team members
\newcommand{\memberone}{Bc. Akhmadinurov Amal}
\newcommand{\membertwo}{Bc. Danilov Denis}
\newcommand{\memberthree}{Bc. Dmitriev Anton}
\newcommand{\memberfour}{Bc. Gainetdinov Ruslan}
\newcommand{\memberfive}{Bc. Kuzin Timofei}
\newcommand{\membersix}{Bc. Tikhoblazhenko Alena}
\newcommand{\memberseven}{Bc. Zimnitski Aliaksei}

% Sprint information -- FILL IN
\newcommand{\sprintname}{Sprint 0}
\newcommand{\sprintperiod}{15 October -- 25 October 2026}  % planned until 29 Oct, closed on 25 Oct
\newcommand{\statusdate}{25 October 2026}                   % date the statuses refer to


% ==================================================
% 2. PACKAGES
% ==================================================

\usepackage[
    a4paper,
    top=3cm,
    bottom=2.7cm,
    left=2.8cm,
    right=2.8cm,
    headheight=45pt,
    headsep=20pt
]{geometry}

\usepackage[T1]{fontenc}
\usepackage[utf8]{inputenc}
\usepackage[english]{babel}
\usepackage{lmodern}

\usepackage{graphicx}
\usepackage[table]{xcolor}
\usepackage{fancyhdr}
\usepackage{titlesec}
\usepackage{setspace}
\usepackage{lastpage}
\usepackage{array}
\usepackage{float}
\usepackage{caption}
\usepackage{booktabs}
\usepackage{tabularx}
\usepackage{csquotes}

% No bibliography in progress reports, so biblatex is not loaded.

\usepackage[hidelinks]{hyperref}


% ==================================================
% 3. COLORS
% ==================================================

\definecolor{primary}{HTML}{000000}
\definecolor{secondary}{HTML}{000000}
\definecolor{muted}{HTML}{000000}
\definecolor{linecolor}{HTML}{000000}
\definecolor{naomiTableHeader}{HTML}{17445A}
\definecolor{naomiTableStripe}{HTML}{EDF7FA}
\definecolor{naomiTableRule}{HTML}{66BBD0}

\newcolumntype{L}[1]{>{\raggedright\arraybackslash}p{#1}}
\newcolumntype{Y}{>{\raggedright\arraybackslash}X}


% ==================================================
% 4. TYPOGRAPHY
% ==================================================

\setstretch{1.2}
\setlength{\parindent}{0pt}
\setlength{\parskip}{7pt}

\titleformat{\section}{\Large\bfseries\color{primary}}{\thesection}{0.8em}{}
\titleformat{\subsection}{\large\bfseries\color{primary}}{\thesubsection}{0.8em}{}
\titleformat{\subsubsection}{\normalsize\bfseries\color{primary}}{\thesubsubsection}{0.8em}{}

\titlespacing*{\section}{0pt}{20pt}{10pt}
\titlespacing*{\subsection}{0pt}{14pt}{7pt}
\titlespacing*{\subsubsection}{0pt}{12pt}{5pt}

\captionsetup{font=small, labelfont=bf, justification=centering}


% ==================================================
% 5. HEADER AND FOOTER
% ==================================================

\newcommand{\naomiheaderfooter}{%
    \fancyhf{}
    \fancyhead[L]{\includegraphics[width=4cm,height=1.15cm,keepaspectratio]{assets/horizontal.png}}
    \fancyhead[R]{%
        \begin{tabular}{r}
            {\footnotesize\bfseries \worksubtitle}\\[3pt]
            {\scriptsize Academic Year \academicyear}
        \end{tabular}}
    \fancyfoot[L]{{\scriptsize \teamname\enspace|\enspace\subject}}
    \fancyfoot[R]{{\scriptsize \thepage\enspace/\enspace\pageref{LastPage}}}
    \renewcommand{\headrulewidth}{0.7pt}
    \renewcommand{\footrulewidth}{0.4pt}%
}

\pagestyle{fancy}
\naomiheaderfooter
\fancypagestyle{plain}{\naomiheaderfooter}


% ==================================================
% 6. TITLE PAGE
% ==================================================

\newcommand{\makemytitlepage}{
    \begin{titlepage}
        \thispagestyle{empty}
        \centering

        % University logo
        \vspace*{0.3cm}
        \includegraphics[width=6cm,height=5cm,keepaspectratio]{assets/vertical.png}

        \vspace{0.5cm}
        {\large\bfseries \university\par}
        \vspace{0.2cm}
        {\normalsize \faculty\par}

        % Project and document title
        \vspace{1.7cm}
        {\small\MakeUppercase{\worksubtitle}\par}
        \vspace{0.5cm}
        \begin{minipage}{0.95\textwidth}
            \centering
            {\fontsize{25}{32}\selectfont\bfseries \worktitle\par}
        \end{minipage}

        % Sprint period and author
        \vspace{0.6cm}
        {\normalsize \sprintperiod\par}

        % Team information
        \vfill
        \begin{minipage}{0.9\textwidth}
            \centering
            \begin{minipage}[t]{0.46\linewidth}
                \centering
                {\normalsize\bfseries Team Leaders\par}
                \vspace{0.25cm}
                {\small
                \leaderone\par
                \leadertwo\par}
            \end{minipage}%
            \hfill
            \begin{minipage}[t]{0.48\linewidth}
                \centering
                {\normalsize\bfseries Team Members\par}
                \vspace{0.25cm}
                {\small
                \memberone\par
                \membertwo\par
                \memberthree\par
                \memberfour\par
                \memberfive\par
                \membersix\par
                \memberseven\par}
            \end{minipage}
        \end{minipage}

        % Bottom information
        \vfill
        {\small \location\enspace\textbullet\enspace\academicyear\par}
    \end{titlepage}
}


% ==================================================
% 7. IMAGES
% ==================================================

% \imageplaceholder{Text}{Caption} -- until the screenshot is added
\newcommand{\imageplaceholder}[2]{
    \begin{figure}[H]
        \centering
        \begingroup
        \setlength{\fboxsep}{0pt}
        \fbox{
            \begin{minipage}[c][5cm][c]{0.85\textwidth}
                \centering
                {\Large\bfseries IMAGE}
                \vspace{0.3cm}

                {\normalsize #1}
            \end{minipage}
        }
        \endgroup
        \caption{#2}
    \end{figure}
}

% \documentimage{path}{caption}{label}
\newcommand{\documentimage}[3]{
    \begin{figure}[H]
        \centering
        \includegraphics[width=0.85\textwidth,keepaspectratio]{#1}
        \caption{#2}
        \label{#3}
    \end{figure}
}


% ==================================================
% 8. TABLES
% ==================================================

\newenvironment{naomitable}[2]{%
    \begin{table}[H]
        \centering
        \caption{#1}
        \label{#2}
        \begingroup
        \small
        \renewcommand{\arraystretch}{1.4}
        \setlength{\tabcolsep}{6pt}
        \rowcolors{2}{naomiTableStripe}{white}
}{%
        \endgroup
    \end{table}
}


% ==================================================
% 9. DOCUMENT
% ==================================================

\begin{document}

\ifshowtitlepage
    \makemytitlepage
\fi


% ==================================================
% SUMMARY
% ==================================================

\section{Summary}

During \sprintname{} (\sprintperiod), the \teamname{} team completed
\textbf{16 of 17 tasks}. The sprint was planned until 29 October and closed
on 25 October, after all but one task had been finished.
The tasks belonged to two epics, \emph{Project Management} and
\emph{Development Environment Setup}, which continue in the next sprints.
The work was distributed among all \textbf{7 team members}.
\sprintname{} was a preparatory sprint: the team set up its tools
(Jira, GitHub organization, documentation repository), agreed on roles and
methodologies, and prepared access to the existing solution and the lab.
Task results were reviewed according to the team's
\emph{Definition of Done} methodology.


% ==================================================
% TASK OVERVIEW
% ==================================================

\section{Task Overview}

% Source: Jira backlog of SCRUM Sprint 0 at sprint closure.

\begin{naomitable}{Tasks of \sprintname{} (status as of \statusdate)}{tab:tasks}
    \footnotesize\renewcommand{\arraystretch}{1.25}
    \begin{tabularx}{\textwidth}{>{\raggedright\arraybackslash\mbox\bgroup}p{1.75cm}<{\egroup} Y L{2.9cm} L{1.9cm} L{1.1cm}}
        \rowcolor{naomiTableHeader}
        \textcolor{white}{\bfseries Key} &
        \textcolor{white}{\bfseries Task} &
        \textcolor{white}{\bfseries Assignee} &
        \textcolor{white}{\bfseries Status} &
        \textcolor{white}{\bfseries Sprint} \\
        SCRUM-13 & Create team website base & A. Dmitriev & Done & 0 \\
        SCRUM-14 & Set up Jira project & R. Gainetdinov & Done & 0 \\
        SCRUM-15 & Assign team roles & R. Gainetdinov & Done & 0 \\
        SCRUM-16 & Discuss roles and Scrumban with supervisor & A. Akhmadinurov & Done & 0 \\
        SCRUM-17 & Write minutes from first supervisor meeting & T. Kuzin & Done & 0 \\
        SCRUM-18 & Write project description and goals & A. Zimnitski & Done & 0 \\
        SCRUM-19 & Create GitHub organization & A. Tikhoblazhenko & Done & 0 \\
        SCRUM-20 & Fork existing repository & A. Tikhoblazhenko & Done & 0 \\
        SCRUM-21 & Arrange lab keys & D. Danilov & Done & 0 \\
        SCRUM-22 & Connect GitHub to Jira & A. Tikhoblazhenko & Done & 0 \\
        SCRUM-24 & Ask our Product Owner about project details in meeting 23.10 & D. Danilov & Done & 0 \\
        SCRUM-25 & Create documentation methodology & A. Tikhoblazhenko & Done & 0 \\
        SCRUM-26 & Create technical methodology & A. Zimnitski & Done & 0 \\
        SCRUM-28 & Request and secure faculty server & T. Kuzin & In Progress & 0, 1 \\
        SCRUM-29 & Replace Team Website epic with Human-Robot Interaction epic & R. Gainetdinov & Done & 0 \\
        SCRUM-31 & Create repo for the team documents & A. Zimnitski & Done & 0 \\
        SCRUM-42 & Set up branch protection and merge request rules & A. Tikhoblazhenko & Done & 0 \\
        \arrayrulecolor{naomiTableRule}\bottomrule[1pt]
    \end{tabularx}
\end{naomitable}


% ==================================================
% TEAM PERFORMANCE
% ==================================================

\section{Team Performance}

The team did not use story points in \sprintname{}, so the velocity report
shows no values for this sprint. Story points are used from Sprint 1, and
velocity will become meaningful after several sprints. Team performance in
\sprintname{} is therefore shown with the sprint burndown chart, which
counts tasks.

\documentimage
    {burndown.png}
    {Sprint burndown chart of \sprintname{} (number of remaining tasks)}
    {fig:burndown}

The chart shows the number of remaining tasks in the sprint.
Rises in the line are tasks added during the sprint: the initial set on
15 October, a larger batch after the team meeting on 20--21 October and a
few tasks after the supervisor meeting on 23 October.
Each batch was completed within one to three days.
At the end of the sprint one task remained open (SCRUM-28).


% ==================================================
% CUMULATIVE FLOW
% ==================================================

\section{Cumulative Flow}

The cumulative flow diagram covers \sprintname{} and the start of Sprint~1
(15--31 October). It counts all work items of the project, including epics.

\documentimage
    {cumulative-flow.png}
    {Cumulative flow diagram for \sprintname{} and the start of Sprint 1}
    {fig:cfd}

The green area (Done) shows completed work items, which grew steadily to
about 30 by the end of October. The purple area (To Do) shows work items
waiting to be done; its jumps on 21, 23 and 29--30 October correspond to new
tasks created after team and supervisor meetings and to planning of Sprint~1.
The bands for In Progress (cyan) and In Review (blue) stay narrow, which means
tasks moved to Done quickly and no stage became a bottleneck. The Testing
status (yellow) was almost unused, because most tasks in \sprintname{} were
organizational and did not involve code.
Most epics remain open, because they cover work planned for the following
sprints.


\end{document}